\section{Conditional dense targets and small-kernel counts}
\label{p14:comp:transfer}

These are the two geometric proofs used in
Section~\ref{p13:transfer:section}. Their source counts, widths, and rank
thresholds are fixed before $N$ tends to infinity. They do not use the
finite-transfer theorem or the law approximation.

\subsection{Proof of Lemma~\ref{p13:transfer:dense-targets}}
\label{p14:comp:transfer-dense-targets}

\begin{proof}
We first count target pairs, keeping the number of binary components in
\emph{each} target explicit. All component counts below include isolated
active coordinates.

\paragraph{The two injective projections.}
Put
\[
 R=\ker\bigl(W\oplus S_y\oplus S_z\longrightarrow W+S_y+S_z\bigr),
 \qquad \dim R=d.
\]
If an element of $R$ has zero $y$-projection, its other two entries
belong to $W\cap S_z=0$. Hence the projection $R\to S_y$ is
injective; the other projection is injective by $W\cap S_y=0$.
In particular $d\leq\min(\dim S_y,\dim S_z)\leq L$.
Represent $W$ by a linear section in the physical source coordinates,
and each target space by a section in its formal even-selection space.
These sections preserve both injective target projections and all actual
prime-valuation equations, including on an irregular window.

Fix $0<\delta<1/2$, to be chosen inside the final exponent margin,
and choose fixed positive $\xi,\theta$ with
\[
 \xi\leq\tfrac14\min(\eta,\delta),\qquad
 \theta\leq\tfrac18\min(\eta,\delta).
\]
Apply Lemma~\ref{p1:fixed:recovery} to obtain $R'\subseteq R$ with
$d':=\dim R'\geq d-\xi L$. Choose its recovery parameter large
enough to give this loss and to control circuits in the $s+2$ physical
blocks. Every nonzero coordinate label on $R'$ then has a fixed
bounded multiplicity $K$.

For each prime $\ell>H\asymp B$, pair equal active labels in its
row. All edges of the resulting binary graph join equal nonzero
labels, so every component lies within one label class and has at
most $K$ vertices. The coordinate labels in either target span
$(R')^*$, since its projection is injective. A component contributes
only one label to this span, even if it meets that target several times.
Consequently each target meets at least $d'$ binary components.
This is a separate lower bound for each target, not just for their union.

\paragraph{Enumerating coefficients from the fixed sources.}
After the equal labels have been paired, the remaining distinct labels
in a prime row split into minimal circuits. Each such circuit uses at
least three physical blocks: $\ell>H$ has at most one occurrence in
any block. Since only two blocks are free, every circuit contains a
source occurrence. The product of primes assigned to a circuit support
therefore divides a fixed source value.
The recovery bound gives a fixed bound on the number of circuit supports
incident to any selected coordinate. For a fixed number of selected
binary components, enumerate their positions, these circuit supports,
one source occurrence on each support, and divisors of the corresponding
source values. The cost is
$B^{O(1)}N^{o(1)}=N^{o(1)}$. Small-prime parities cost
$2^{O(\pi(H))}=N^{o(1)}$ as well. Thus the squarefree coefficients
of their component equations have $N^{o(1)}$ choices and height
$N^{A_0}$ for a fixed $A_0$. The height bound per component is fixed
once $s$, the recovery parameters and the width multiplier are fixed;
it does not increase when more disjoint component equations are selected.
The number of equations required by Lemma~\ref{p1:arith:geometric}
is chosen afterwards. No subspace or full deletion pattern is enumerated.

\paragraph{Removing the small aligned curves.}
Two exact units between the targets determine a reduced alignment of
height $O(B)$. Two distinct alignments of this height would force
$y,z=O(B^3)$, impossible in the critical range. If there are fewer than
two exact units, no removal is needed. Otherwise let $u\geq2$ be the
number of units on the unique alignment. Their even selections have dimension
$u-1$ before passage to square classes, and a target has at most one
local rank defect. Their common image therefore gives
\[
 u-2\leq\dim(S_y\cap S_z)\leq\theta L
\]
in the weak aligned case. Keep at most one of these units and exclude
from selection every binary component meeting an endpoint of another
unit. At most $2(u-1)\leq2\theta L+2$ components are excluded.
This deletes components from the selectable collection; it does not
restrict $R'$ again. The surviving positions contain no two exact units
on an alignment, which excludes the small-curve alternatives in the
proof of Lemma~\ref{p1:arith:geometric}. No removal is needed for a
nonaligned pair.

For $t=y,z$, write $m_t$ for the number of selectable components
meeting target $t$. The separate projection bounds now give
\begin{equation*}
 m_t\geq d'-2\theta L-2
      \geq d-(\xi+2\theta)L-2.
\end{equation*}
Thus $m_t\geq\eta L/3$ for large $N$ whenever $d\geq\eta L$;
if $d>(1/2+\delta)L$, then
$m_t>(1/2+\delta/3)L$ for both targets.

\paragraph{The first stratum: a positive rank proportion.}
A component type records its number of coordinates in each physical
block. There are finitely many types, because the component size is at
most $K$. If a type meeting both targets occurs sufficiently often,
select the fixed number of its components required by the two-variable
geometric count. With the small curves excluded above and coefficients
of height $N^{A_0}$, the proof of Lemma~\ref{p1:arith:geometric}
gives $N^{1/2+\varepsilon'}$ target pairs, including the
$N^{o(1)}$ pattern and coefficient choices.

Otherwise only a bounded number of components meet both targets.
Each target still has a positive proportion of $L$ components avoiding
the other target. In each target choose a univariate type occurring with
positive frequency, and then the fixed number of its components needed
by the one-variable part of Lemma~\ref{p1:arith:geometric}, taking
at least two. Their component equations form a one-variable system
with fixed coefficients and at least two distinct selected offsets;
different components have disjoint vertices. In the degree-one case,
these are precisely the two linear factors for
Lemma~\ref{p1:arith:pell}. The one-variable count gives
$N^{\varepsilon'}$ target pairs after allocating the margin between
the two centers and the coefficient choices. Every selected set has fixed size, so its even-selection
square-class map is injective for large $N$ by
Lemma~\ref{p1:arith:runge}; irregularity of the entire window is
irrelevant to this local use. We have proved
\[
 \#\{(y,z):d\geq\eta L\}\ll N^{1/2+\varepsilon'}.
\]
Both counts here and below impose the pair restriction in the lemma.

\paragraph{The second stratum: rank greater than half a window.}
Suppose $d>(1/2+\delta)L$. Let $n_{1,t}$ be the number of selectable
components using exactly one coordinate in target $t$. Every other
component meeting that target uses at least two of its $B=L+1$
coordinates. Hence, separately for $t=y,z$,
\[
 B\geq n_{1,t}+2(m_t-n_{1,t}),\qquad
 n_{1,t}\geq2m_t-B>\frac{2\delta}{3}L-1.
\]
If a mixed type with one coordinate in either target occurs sufficiently
often, the degree-one case of the geometric count gives
$N^{\varepsilon'}$ pairs directly. Otherwise only boundedly many of
the $n_{1,t}$ components are mixed. The displayed inequality then
provides two components for each target, each using one coordinate of
that target and no coordinate of the other. The resulting two distinct
linear factors give the uniform two-value Pell count for each center.
Thus
\[
 \#\{(y,z):d>(1/2+\delta)L\}\ll N^{\varepsilon'}.
\]
This explains the threshold $L/2$: above it, the coordinate budget
forces singleton occurrences in each target.

Finally $p_y,p_z\ll N^{-1}$ and $2^L\asymp_C N$.
The stratum $\eta L\leq d\leq(1/2+\delta)L$ has weighted cost
$N^{1/2+\varepsilon'}N^{-2+1/2+\delta}$, and the stratum above it
has cost $N^{\varepsilon'}N^{-2+1}$. Choose $\delta>0$ and then
$\varepsilon'>0$ so that $\delta+\varepsilon'<\varepsilon$;
all preceding choices are fixed before $N$. This proves
\eqref{p13:transfer:dense-target-bound}.
In a primitive family, each target is direct from the prescribed
$r-2$ sources separately, as well as from the other target; the full
nullity is at most $L$. Multiplication by the $r-2$ source marginals
and summation over the prescribed source collection gives
\eqref{p13:transfer:restricted-dense}, after dividing the final exponent
margin between the source count and the conditional estimate.
\end{proof}

\subsection{Proof of Lemma~\ref{p13:transfer:small-kernels}}
\label{p14:comp:transfer-small-kernels}

\begin{proof}
For the first assertion fix the $m$ offsets, at a polylogarithmic cost,
and write $d_i=n_i/K_Y(n_i)=s_i a_i^2$. Since
$d_i\gg N^{1-b}$ and $\gcd(d_i,d_j)\ll B$,
\[
 \operatorname{lcm}(d_1,\ldots,d_m)
 \geq\frac{\prod_i d_i}{\prod_{i<j}\gcd(d_i,d_j)}
 \geq N^{m(1-b)-o(1)}>N.
\]
The $a_i$ are pairwise coprime. If $\prod_i a_i^2>N$, multiply the
square factors in a fixed order until their product first exceeds $N$.
The resulting square $D$ lies between $N$ and $N^{1+u}$, so it has
at most $N^{(1+u)/2}$ possible roots. Assign its complete prime powers
to the $m$ offsets; the number of assignments is
$m^{\omega(D)}=N^{o(1)}$. The CRT leaves at most one center modulo
$D>N$ in the interval.

If $\prod_i a_i^2\leq N$, start with this product and add the small
prime factors from the least common multiple until the first crossing
of $N$. Now $N<D\leq NY$ and $D=sa^2$ with $s$ $Y$-smooth.
For fixed $\sigma>\max(1/2,1-1/A)$, the number of such modules is at most
\[
 (NY)^\sigma\sum_{a\geq1}a^{-2\sigma}
       \prod_{\ell\leq Y}(1-\ell^{-\sigma})^{-1}
       =N^{\sigma+o(1)}.
\]
Choose $\sigma$ just above $1-1/A$. Assigning the complete prime powers
again costs $N^{o(1)}$ and the CRT gives at most one center. This proves
\eqref{p13:transfer:small-kernels-bound}.

For the second assertion, use $H=C_DB$ as in
Lemma~\ref{p13:half:small-kernels}.
That lemma gives the result for $q>0$.
For $q=0$, use $\mathcal X_0\subseteq\mathcal X_{\varepsilon/2}$
and apply the lemma with error margin $\varepsilon/2$.
\end{proof}
